% =============================================================================
% Response to Reviewers -- Journal of Air Transport Management
% Manuscript: A Kirkpatrick-based Framework for Integrating Safety and
%             Sustainability in Aviation Maintenance: A Case Study on Waste
%             Reduction and Resource Efficiency
% =============================================================================
\documentclass[11pt,a4paper]{article}
\usepackage[utf8]{inputenc}
\usepackage{geometry}\geometry{margin=2.5cm}
\usepackage[expansion=false]{microtype}
\usepackage{parskip}
\usepackage{xcolor}
\usepackage{enumitem}
\usepackage{hyperref}
\hypersetup{colorlinks=true, linkcolor=blue!60!black, urlcolor=blue!60!black}

\newcommand{\comment}[1]{\par\smallskip\noindent\textbf{\color{blue!60!black}Reviewer comment:}\\ \emph{#1}\par\smallskip}
\newcommand{\response}[1]{\noindent\textbf{Response:} #1\par\smallskip}
\newcommand{\change}[1]{\par\noindent\textbf{\color{green!50!black}Change in manuscript:}\\ \emph{#1}\par\smallskip}

\title{Response to Reviewers}
\author{\vspace{-1em}}
\date{}

\begin{document}
\maketitle

\noindent\textbf{Manuscript:} A Kirkpatrick-based Framework for Integrating
Safety and Sustainability in Aviation Maintenance: A Case Study on Waste
Reduction and Resource Efficiency \\
\noindent\textbf{Journal:} Journal of Air Transport Management \\
\noindent\textbf{Submission type:} Case Study (revised)

\section*{Letter to the Editor}

Dear Editor,

We sincerely thank you and the Reviewer(s) for the detailed and constructive
feedback provided on our manuscript. We have carefully addressed every
comment, and the revised version incorporates substantial structural,
analytical and presentational changes. Below we respond point-by-point to
each of the ten comments received. Where applicable, we cite the section,
table or figure in the revised manuscript where the change can be found.

We believe the revised paper now responds fully to the editorial focus of
the \emph{Journal of Air Transport Management} on managerial relevance,
policy implications and quantitative outcomes.

\section*{Point-by-point response}

% -----------------------------------------------------------------------------
\subsection*{Comment 1 -- Add a Dedicated ``Managerial Implications'' Section}
\comment{The manuscript would benefit from a dedicated Managerial Implications
section covering cost-benefit justification, integration with existing SMS
and ERP infrastructure, audit-readiness, and sustainability/ESG reporting.}

\response{We agree and have added a new Section~6 \emph{Managerial
Implications}, structured exactly along the four sub-themes suggested
(cost-benefit, SMS/ERP integration, audit-readiness against AS9100~§7.2,
sustainability/ESG reporting). The cost-benefit subsection presents the
ROI calculation and sensitivity bounds. The integration subsection provides
a five-step roll-out checklist designed for non-disruptive integration. The
audit-readiness subsection maps each Kirkpatrick level to specific AS9100
clauses. The sustainability subsection lists four auditable ESG indicators.}

\change{New Section~6 (Managerial Implications) with subsections 6.1--6.4.}

% -----------------------------------------------------------------------------
\subsection*{Comment 2 -- Add a Short ``Policy Recommendations'' Subsection}
\comment{A short Policy Recommendations subsection (after Conclusion or as
part of Discussion) addressed to FAA/EASA/ANAC, IATA and SAE would
significantly strengthen the alignment of the paper with the editorial focus
of JATM.}

\response{We added a dedicated Section~7 \emph{Policy Recommendations} with
three numbered recommendations: (1)~embed training-effectiveness evaluation
in SMS acceptance criteria; (2)~develop joint IATA/SAE industry guidance on
training evaluation for MRO; (3)~mandate annual reporting of
training-effectiveness metrics as safety indicators. We also include a
paragraph on coupling with sustainability disclosure regimes (e.g., ESRS).}

\change{New Section~7 (Policy Recommendations) with subsections 7.1--7.3.}

% -----------------------------------------------------------------------------
\subsection*{Comment 3 -- Strengthen Literature Review with JATM-Published Papers}
\comment{Adding 3--5 recent JATM articles (2023--2025) on maintenance
training, MRO sustainability, and human factors would significantly improve
the alignment of the paper with the journal scope.}

\response{We expanded the literature review (Section~2) to cite five
additional recent JATM papers on the compliance paradox in aviation training
\citep{Nakamura2024}, regulatory gaps in Kirkpatrick implementation
\citep{ORreilly2024}, digital twins in MRO training \citep{Garcia2024},
training-failure costs in MRO benchmarking \citep{Zhang2025}, and predictive
ROI modelling \citep{Peterson2025}. The reference list in the revised
manuscript thus includes papers from JATM volumes 116--123 (2024--2025)
alongside the previously cited classic and methodological references.}

\change{Additions throughout Section~2 (Literature Review); five new entries
added to \texttt{references.bib}.}

% -----------------------------------------------------------------------------
\subsection*{Comment 4 -- Add a Comparative Table of Key Results}
\comment{A Before/After comparative table (training failures,
non-conformances, quality costs, customer satisfaction, material waste)
would be powerful for JATM reviewers.}

\response{We added Table~1 \emph{Summary of operational and sustainability
outcomes before and after the Kirkpatrick-based intervention}. The table
reports seven metrics across Before/After/Variation columns: training
failures, regulatory non-conformances, first-attempt Level~2 pass rate,
Level~3 work-card consultation rate, quality costs, turnaround time and
customer satisfaction. Pre-intervention values are anchored to the
facility's trailing twelve-month records, while post-intervention values
cover the twelve-month intervention period.}

\change{New Table~1 referenced at the end of Section~4.5 (Results).}

% -----------------------------------------------------------------------------
\subsection*{Comment 5 -- Elaborate on the ``Economic Analysis'' Subsection}
\comment{The economic analysis in the original Section~4.4 is descriptive
rather than quantitative. A cost-of-quality decomposition (internal/external
failure, prevention/appraisal) and an explicit ROI calculation would
strengthen the paper substantially.}

\response{Section~4.6 \emph{Economic analysis: impact on cost of quality}
now presents a full CoQ decomposition. We provide formulas for internal
failure cost reduction $\Delta C_{\text{IF}}$ (rework labour hours $\times$
loaded rate plus scrapped components $\times$ mean material cost), external
failure cost reduction $\Delta C_{\text{EF}}$ (warranty claims plus
expedited shipping plus customer concessions), and prevention/appraisal
costs $C_{\text{P\&A}}$. We then compute the ROI explicitly using
Eq.~(3) and report sensitivity bounds ($\pm 25\%$ on all reported
reductions). Even under the conservative scenario the ROI remains in the
high triple digits, supporting the robustness of the economic case.}

\change{Expanded Section~4.6 with equations (1)--(3) and sensitivity
analysis.}

% -----------------------------------------------------------------------------
\subsection*{Comment 6 -- Add a ``Limitations and Future Research'' Table}
\comment{The current limitations are presented as prose. A structured table
(Limitation $\to$ Mitigation $\to$ Future Research) would improve clarity
and signal intellectual honesty.}

\response{We have converted the limitations from prose into a structured
Table~2 with three columns (Limitation / Mitigation in the present study /
Future research direction). Six limitations are catalogued: single-case
design, inability to isolate training effects, self-reported behaviour
data, generalisability to smaller MROs, order-of-magnitude sustainability
outcomes, and aggregated-data-only access. The body text of
Section~\ref{sec:limitations} also expands on three broader future
research avenues (cross-regulatory comparison, digital-twin integration,
formal causal attribution).}

\change{New Table~2 in Section~9 (Limitations and future research) plus
expanded narrative.}

% -----------------------------------------------------------------------------
\subsection*{Comment 7 -- Revise the Abstract for JATM Style}
\comment{The abstract should lead with economic outcomes and concrete
numbers. Current abstract is more theoretical in tone.}

\response{We rewrote the abstract so that the economic outcomes appear in
the first third (after the problem statement): the 28\% reduction in
training failures, 35\% reduction in regulatory non-conformances, the
USD~1.2\,M annual quality-cost savings, and the customer-satisfaction
improvement from 3.4 to 4.7. The final third explicitly states the policy
implication for ICAO Annex~19. The revised abstract is approximately 270
words.}

\change{Abstract rewritten in \texttt{main.tex} frontmatter.}

% -----------------------------------------------------------------------------
\subsection*{Comment 8 -- Add a Graphical Conceptual Model (Figure 0 or in Introduction)}
\comment{A new conceptual figure mapping Kirkpatrick levels to AS9100
clauses, SMS pillars and sustainability outcomes would strengthen the
visual narrative and could double as the graphical abstract.}

\response{We added Figure~1 \emph{Conceptual model linking Kirkpatrick's
four levels to AS9100 clauses, SMS pillars and sustainability outcomes},
placed at the end of the Introduction. The figure shows four rows (one per
Kirkpatrick level) and three columns (Kirkpatrick, AS9100/SMS,
Sustainability), with horizontal arrows representing the mapping proposed
in this study and vertical dashed arrows representing the progression of
the Kirkpatrick model. A condensed three-block version is used in the
Graphical Abstract.}

\change{New Figure~1 (\texttt{figs/fig0-conceptual-model.tex}) included at
the end of Section~1.2; Graphical Abstract uses condensed version.}

% -----------------------------------------------------------------------------
\subsection*{Comment 9 -- Clarify ``Training Failures''}
\comment{An operational definition of \emph{training failure} would improve
methodological transparency.}

\response{We added an explicit operational definition in
Section~3.2: ``A \emph{training failure} is an instance where a maintenance
technician, after completing a scheduled training intervention, commits a
documented procedural error within the subsequent 90-day window that is
directly attributable to a knowledge or skill gap identified in the
post-training evaluation. Errors caused by tool malfunction, supplier
non-conformance or documented changes in the standard procedure are
excluded.''}

\change{New definition block in Section~3.2 (Data collection).}

% -----------------------------------------------------------------------------
\subsection*{Comment 10 -- Check and Update References to JATM Format}
\comment{Ensure all references follow JATM/Elsevier formatting (author-year,
journal names in italics, DOIs where available).}

\response{We have audited the bibliography and rebuilt it in BibTeX format
using \texttt{cas-model2-names.bst} (Elsevier author-year). Every entry now
follows the same field structure (author / title / journal / volume /
pages / year / DOI). Entries for which a specific field could not be
verified against the original source carry an explicit \texttt{\% TODO}
flag in the \texttt{references.bib} source. We will resolve the remaining
flags prior to the next submission round.}

\change{Bibliography reformatted in \texttt{references.bib}; compilation
now uses \texttt{cas-model2-names} BibTeX style.}

\section*{Closing}

We are grateful for the time and care invested by the Editor and the
Reviewer(s) and believe that the revisions described above respond fully
to the feedback received. We remain available for any additional revision
that may be requested.

\end{document}
